


Soit la base de données suivante permettant de gérer un championnat de football.
\begin{itemize}
\item Stade(\textbf{id}, ville, nom, nbPlaces, prixBillet)
\item Equipe(\textbf{id}, pays, siteWeb, entraîneur)
\item Joueur(\textbf{id}, \textit{idEquipe}, nom, prénom, âge)
\item Match(\textbf{id}, \textit{idStade}, dateMatch, \textit{idEquipe1}, \textit{idEquipe2}, 
      scoreEquipe1, scoreEquipe2, nbBilletsVendus)
\item But(\textbf{\textit{idJoueur}, \textit{idMatch}, minute}, penalty)

\end{itemize}


Tous les attributs \textbf{en gras}
sont des clés primaires, ceux \textit{en italiques} sont des clés étrangères. Notez
qu'une clé étrangère peut faire partie d'une clé primaire.

Chaque pays a une seule équipe en compétition. Une ville peut avoir plusieurs stades. L'attribut But(penalty) vaut VRAI si le but a été marqué suite a un penalty (coup de pied de réparation) ou FAUX dans le cas contraire. Le nom de l'entraîneur d'une équipe est donné par l'attribut Equipe(entraîneur). 

\section{Conception (6 points)}

Etudions la conception de cette base en répondant aux questions suivantes.

\begin{enumerate}
%   \item Avec ce schéma, une équipe peut-elle jouer contre elle-même\,? Expliquez.
    
    \item Avec ce schéma, un joueur peut-il marquer un but dans un match auquel
    son équipe ne participe pas\,? Expliquez.

    \item Est-il possible de savoir dans quelle ville a été marqué chaque but\,? Expliquez.
    
     \item Donnez le schéma entité-association correspondant aux relations \textsc{Match}, \textsc{Equipe},
          \textsc{Stade}.
          
     \item Donnez la commande SQL de création de la table \textsc{Match}.
     
     \item Citer au moins une clé \emph{candidate} autre que la clé primaire 
     parmi les attributs des tables (bien lire l'énoncé). 
       Quel serait l'impact si on la choisissait comme clé primaire\,?
       
     \item Un de vos collègues propose de modéliser un but comme une association plusieurs-plusieurs
     entre un joueur et un match. Cela vous semble-t-il une bonne idée\,? Cela correspond-il 
     au schéma relationnel de notre base\,?
\end{enumerate}

\begin{solution}
\begin{enumerate}
  \item Oui, la base permet de représenter un but marqué dans un match opposant deux équipes
       dont aucune n'est celle du joueur.
  \item Oui, par transitivité But $\to$ Match, Match $\to$ Stade  et Stade $\to$ Ville.
  \item Schéma
  \item 

\begin{minted}{sql}
CREATE TABLE Match(
    id NUMBER(10), 
    idStade NUMBER(10), 
    dateMatch DATE, 
    idEquipe1 NUMBER(10), 
    idEquipe2 NUMBER(10), 
    scoreEquipe1 NUMBER(10), 
    scoreEquipe2 NUMBER(10),
    PRIMARY KEY (id),
    FOREIGN KEY (idStade) REFERENCES Stade(id),
    FOREIGN KEY (idEquipe1) REFERENCES Equipe(id),
    FOREIGN KEY (idEquipe2) REFERENCES Equipe(id)
);
\end{minted}
  
  \item Le pays est clé candidate pour l'équipe. Il faudrait alors l'utiliser comme clé étrangère 
     partout.
     
\end{enumerate}

\end{solution}

\section{SQL (8 points)}

Exprimez les requêtes suivantes en SQL.

\begin{enumerate}

\item Noms des joueurs âgés de plus de 30 ans qui ont 
     marqué un but dans la première minute de jeu. 

\begin{solution}
\begin{minted}{sql}
select nom 
from Joueur as j, But as b
where j.id = b.idJoueur and j.âge > 30 and b.minute = 1;
\end{minted}
\end{solution}

\item Noms des joueurs français qui n'ont marqué aucun but

\begin{solution}
\begin{minted}{sql}
select nom 
from Joueur as j, Equipe as e
where j.idEquipe = e.id 
and e.pays = 'France' 
and j.id NOT IN (select idJoueur from But)
\end{minted}
\end{solution}

\item Le prix du billet le plus cher

\begin{solution}

\begin{minted}{sql}
select prixBillet 
from Stade as s1
where not exists (select ''
                from Stade as s2
                where s2.prixBillet >= s1.prixBillet)
\end{minted}
\end{solution}


\item Donnez le nom des entraineurs avec le nombre de buts marqués par leur équipe

\begin{solution}
\begin{minted}{sql}
select entraineur, count(*)
from Equipe as e, Joueur as j, But as b
where e.id=j.idEquipe
and b.idJoueur = j.id
group by e.id, e.entraineur
\end{minted}
\end{solution}



\item Donnez le nom des joueurs ayant marqué un but dans un match auquel leur équipe
    ne participe pas.  

\begin{solution}
\begin{minted}{sql}
select prénom, nom
from Match as m, Joueur as j, But as b
where e.id=j.idEquipe
and b.idJoueur = j.id
and m.id = b.idMatch
and e.id != m.idEquipe1
and e.id != e.idEquipe2
\end{minted}
\end{solution}


\item Les noms des joueurs français qui ont marqué un but lors d'un match entre la France et la Suisse 
    joué à Lille. Attention\,:  dans la table \textsc{Match}
    la France peut être l'équipe 1 ou l'équipe 2. Il faut
      construire la requête qui correspond aux deux cas.
     
\begin{solution}
Il faut soit Match.equipe1='France', soit  Match.equipe2='France'. 
\begin{minted}{sql}
select j.nom
from Match m, Equipe e1, Equipe e2, Joueur j, Stade s, But b
where m.idEquipe1 = e1.id 
and m.idEquipe2 = e2
and ((e1.pays = 'France' id and e2.pays='Suisse' and e1.id = j.idEquipe)
         or
     (e2.pays = 'France' id and e1.pays='Suisse' and e2.id = j.idEquipe))
and m.idStade = s.id 
and s.ville = 'Lille' 
and b.idJoueur = j.id 
and b.idMatch = m.id )
\end{minted}
\end{solution}


\end{enumerate}

\section{Algèbre (3 points)}



\begin{enumerate}
   \item Donnez l'expression algébrique de la première requête SQL (section précédente)

\begin{solution}
$\pi_{nom}(\sigma_{age > 30} (Joueur) \underset{id=idJoueur}{\bowtie} \sigma_{minute = 1}(But) )$
\end{solution}

\item Donnez l'expression algébrique pour la seconde requête (les joueurs français qui ne marquent pas)

\begin{solution}
Joueurs français:
$$A = \sigma_{pays='France'}(Joueur\;j) \underset{j.idEquipe = e.id}{\bowtie} Equipe\;e)$$
Joueurs français qui ont marqué au moins un but :
$$B = \sigma_{pays='France'}(But\;b \underset{b.idJoueur = j.id}{\bowtie} Joueur\;j) \underset{j.idEquipe = e.id}{\bowtie} Equipe\;e)$$

Résultat:\\
$$Resultat := A - B$$
\end{solution}

\item Donnez la requête SQL correspondant à l'expression algébrique suivante et expliquez le sens de cette requête

$$\pi_{idMatch}(But \Join_{idMatch=id} \sigma_{scoreEquipe1=0 \land scoreEquipe2=0}(Match))$$
\begin{solution}
Les matchs nuls 0-0 pour lequel il existe quand même un but.
\begin{minted}{sql}
select idMatch
from But as b, Match as m
where b.idMatch=m.id
and m.scoreEquipe1=0 
and m.scoreEquipe2=0
\end{minted}
\end{solution}
\end{enumerate}

\section{Programmation et transactions (3 points)}

\begin{enumerate}

\item   Parmi les phrases suivantes, lesquelles vous semblent exprimer une contrainte de \emph{cohérence} correcte sur la base de données?

\begin{enumerate}
  \item Pour chaque match nul, je dois trouver autant de lignes associées au match dans la table \texttt{But} pour l'équipe 1 et pour l'équipe 2.
  \item Quand j'insère une ligne dans la table \texttt{But} pour l'équipe 1, je dois effectuer également une insertion pour l'équipe 2 dans la table \texttt{But}.
  \item Si je trouve une ligne dans la table \texttt{Match} avec \texttt{id=x}, \texttt{scoreEquipe1=y} et \texttt{scoreEquipe2=z}, alors je dois trouver \texttt{y+z} lignes dans la table \texttt{But} avec \texttt{idMatch=x}.
 \end{enumerate}
 
 \item Voici un programme en pseudo-code simplifié à exécuter chaque fois qu'un but est marqué dans un match, par un joueur et au bénéfice de la première équipe. Les '\#' marquent des commentaires.
        
\begin{verbatim}
function marquer_but_eq1 (imatch, ijoueur, min, peno)
     startTransaction
     # A
     select scoreEquipe1 into :scoreE where idMatch = :imatch
     #B
     update Match set scoreEquipe1 = :scoreE + 1 where idMatch = :imatch
     # C
     insert into But (idMatch, idJoueur, minute, penalty)
     values (:idMatch, :ijoueur, :min, :peno)
     # D
\end{verbatim}
 
 Où faut-il selon-vous ajouter un ordre \texttt{commit}?
 \begin{enumerate}
   \item Juste après \texttt{B} et \texttt{C}
   \item Juste après \texttt{A}
   \item Juste après \texttt{D}
 \end{enumerate}
  Justifiez votre réponse.
  
  \item Supposons que je lance deux exécutions de  \texttt{marquer\_but\_eq1} en même temps.
  Dans quel scénario l'exécution imbriquée amène-t-elle une incohérence\,?
\end{enumerate}
